\documentclass[11pt]{letter}
\usepackage[margin=1in]{geometry}
\usepackage{hyperref}

\signature{Joshua Clingo, PhD and Jeff Yoshimi, PhD \\ School of Social Sciences, Humanities and Arts \\ University of California, Merced \\ jclingo@ucmerced.edu and jyoshimi@ucmerced.edu}

\begin{document}

\begin{letter}{The Editors \\ \emph{Adaptive Behavior} \\ Sage Publications}

\opening{Dear Editors,}

We are pleased to submit our manuscript, \emph{Modeling Decanalization with Homeostatic Reinforcement Learning}, for consideration as an original research article in \emph{Adaptive Behavior}.

The manuscript presents, to our knowledge, the first computational instantiation of the canalization–decanalization–recanalization cycle that has become central theorizing in computational psychiatry, integrative psychotherapy, and work on psychedelic-assisted treatment. Using a homeostatic navigation task, we compared homeostatic reinforcement learning (HRL) against negative-feedback (NF) trained agents under matched canalization and decanalization conditions. HRL agents showed improved post-decanalization homeostatic performance relative to matched no-decanalization controls, while NF agents did not. The result helps demonstrate how temporary destabilization can yield more adaptive behavior.

We believe the work fits the aims and scope of \emph{Adaptive Behavior}: it is an agent-based, embodied simulation that bridges dynamical-systems theorizing, reinforcement learning, and homeostatic regulation, and it renders long-standing metaphors in a form that other modelers can extend or contest. All source code, evaluation data, and analysis materials are openly available at \url{https://github.com/JClingo/decanalization}.

We thank you for considering our submission and look forward to the reviewers' feedback.

\closing{Sincerely,}

\end{letter}

\end{document}
